\documentclass{article}
\usepackage{geometry}
\geometry{margin=1.5cm, vmargin={0pt,1cm}}
\setlength{\topmargin}{-1cm}
\setlength{\headheight}{12.5pt}
\setlength{\paperheight}{29.7cm}
\setlength{\textheight}{25.3cm}

% useful packages.
% \usepackage[UTF8]{ctex}
\usepackage{fancyhdr}
\usepackage{layout}
\usepackage{tikz}
\usetikzlibrary{arrows.meta}

% import common macros
% % % % % % % % % % % % % % % % % % % % % % % % % % % % % % % %
% Common Macros for LaTeX
% % % % % % % % % % % % % % % % % % % % % % % % % % % % % % % %

% Useful packages
\usepackage{amsfonts}
\usepackage{amsmath}
\usepackage{amssymb}
\usepackage{amsthm}
\usepackage{enumerate}
\usepackage{graphicx}
\usepackage{multicol}
\usepackage[ruled,linesnumbered]{algorithm2e}

% Math operators and common symbols
\newcommand{\dif}{\mathrm{d}}
\newcommand{\innerProd}[1]{\left\langle #1 \right\rangle}
\newcommand{\difFrac}[2]{\frac{\dif #1}{\dif #2}}
\newcommand{\pdfFrac}[2]{\frac{\partial #1}{\partial #2}}
\newcommand{\T}{\mathrm{T}}
\newcommand{\norm}[1]{\left\| #1 \right\|}
\newcommand{\abs}[1]{\left| #1 \right|}

% Vector and Matrix shortcuts
\newcommand{\vecTwo}[2]{
  \begin{bmatrix}
    #1 \\
    #2
\end{bmatrix}}
\newcommand{\vecThree}[3]{
  \begin{bmatrix}
    #1 \\
    #2 \\
    #3
\end{bmatrix}}
\newcommand{\vecTwoH}[2]{
  \begin{bmatrix}
    #1 & #2
  \end{bmatrix}
}
\newcommand{\vecThreeH}[3]{
  \begin{bmatrix}
    #1 & #2 & #3
  \end{bmatrix}
}

% Common fields
\newcommand{\R}{\mathbb{R}}
\newcommand{\C}{\mathbb{C}}
\newcommand{\Z}{\mathbb{Z}}
\newcommand{\N}{\mathbb{N}}

% Solutions and proofs
\newenvironment{solution}{
\begin{proof}[解]}{
\end{proof}}

% Utility to get environment variables (requires lualatex)
\newcommand{\getEnv}[2]{\directlua{tex.print(os.getenv("#1") or "#2")}}


% Name and uID
\newcommand{\name}{\getEnv{NAME}{NAME}}
\newcommand{\uid}{\getEnv{UID}{UID}}
\newcommand{\course}{Complex Analysis}
\newcommand{\hwNum}{10}

\begin{document}

\pagestyle{fancy}
\fancyhead{}
\lhead{\zhstr{\name} (\uid)}
\chead{\course\ \#\hwNum}
\rhead{\today}

\section{Exercise 4.31}

Let $l_1, l_2$ be two rays that both start from the origin and have an
angle $\theta \in (0, 2\pi)$. Find the angle between the image curves
$f(l_1)$ and $f(l_2)$ at $0$ by assuming the conditions in
Theorem~4.30.

\begin{solution}
  Consider $m$ as the multiplicity of $f$ at $0$. By Theorem~4.30,
  the function $f$ can be written as:
  \[
    f(z) = [g(z)]^m,\quad g \text{ is conformal}
  \]
  Write $l_1$ and $l_2$ as $l_1 = \{re^{i\alpha} : r > 0\}$ and
  $l_2 = \{re^{i(\alpha + \theta)} : r > 0\}$ for some $\alpha \in
  \R$. Then we have
  \[
    \begin{aligned}
      f(l_1) &= \{[g(re^{i\alpha})]^m : r > 0\} \\
      f(l_2) &= \{[g(re^{i(\alpha + \theta)})]^m : r > 0\}.
    \end{aligned}
  \]
  By the conformality of $g$, the angle between $g(l_1)$ and $g(l_2)$
  at $0$ is $\theta$. Since $f(z) = [g(z)]^m$, the angle between
  $f(l_1)$ and $f(l_2)$ at $0$ is $m\theta$.
\end{solution}

\section{Exercise 4.34}

For a nonconstant analytic function $f : D \to \C$ where $D$ is a
domain, $0 \in D$, $f(0) = f'(0) = 0$, and $f''(0) \ne 0$,
construct two curves $\Gamma_1$ and $\Gamma_2$ so that
\begin{enumerate}[(a)]
  \item both $\Gamma_1$ and $\Gamma_2$ pass through $0$,
  \item $\Gamma_1$ and $\Gamma_2$ are orthogonal at $0$,
  \item $f(\Gamma_1)$ is a subset of $\R$ and has a minimum at $0$,
  \item $f(\Gamma_2)$ is a subset of $\R$ and has a maximum at $0$.
\end{enumerate}

\begin{solution}
  By Theorem~4.30, the function $f$ can be written as
  \[
    f(z) = [g(z)]^2, \quad g: U_r(0) \to \C \text{ is conformal}.
  \]
  Where $g(0) = 0$ and $g'(0) \ne 0$. Thus there exists an inverse function
  $g^{-1} : U_{r'}(0) \to U_r(0)$ for some $r' > 0$. Let $\Gamma_1 =
  g^{-1}(\R \cap U_{r'}(0))$ and $\Gamma_2 = g^{-1}(i\R \cap
  U_{r'}(0))$. Then $\Gamma_1$ and $\Gamma_2$ are orthogonal at $0$
  since $g$ is conformal. Moreover, $f(\Gamma_1) = \{x^2 : x \in \R
  \cap U_{r'}(0)\}$ has a minimum at $0$, and $f(\Gamma_2) = \{-y^2 :
  y \in \R \cap U_{r'}(0)\}$ has a maximum at $0$.
\end{solution}

\section{Exercise 4.55}

Show that there is no analytic function $f$ in $D$ that extends
continuously to $\overline{D}$ and satisfies $f(z) = \bar{z}$ for all
$z \in \partial D$.

\begin{solution}
  If there exists such an analytic function $f$, then we have: $f(z)
  = \bar{z}$ for all $z \in \partial D$. Consider the integral:
  \[
    \oint_{\partial D} f(z) \dd z = \oint_{\partial D} \bar{z} \dd z
  \]
  By Green's theorem, we have
  \[
    \oint_{\partial D} (x-iy)(\dd x + i \dd y) = \iint_D
    \left(\frac{\partial (x-iy)}{\partial x} + i \frac{\partial
    (x-iy)}{\partial y}\right) \dd x \dd y = 2i \iint_D \dd x \dd y =
    2i\text{Area}(D) \ne 0.
  \]
  However, since $f$ is analytic in $D$, we have $\oint_{\partial D}
  f(z) \dd z = 0$. This is a contradiction.
\end{solution}

\section{Exercise 4.57 (Alternative proof of Liouville's theorem)}

Prove Theorem~3.95 by Lemma~4.56.

\begin{solution}
  Consider a bounded entire function $f$, and let $M = \sup_{z \in
  \C} |f(z)| < \infty$. Let $g(z) = f(z) - f(0)$, then $g$ is also a
  bounded entire function with an upper bound of $2M$. Then we define:
  \[
    h_R(x): \{|z| < 1\} \to \C, \quad h_R(x) = \frac{g(Rx)}{2M} =
    \frac{f(Rx) - f(0)}{2M}.
  \]
  Then $h_R$ is an analytic function in $\{|z| < 1\}$ and satisfies
  $|h_R(x)| \le 1$ for all $x$ such that $|x| < 1$. By Lemma~4.56, we have
  \[
    |h_R(x)| \leq |x| \Rightarrow |f(x) - f(0)| \leq \frac{2M}{R} |x|
    \text{ for all } |x| < R.
  \]
  Let $R \to \infty$, we have $|f(x) - f(0)| = 0$ for all $x \in \C$,
  which means $f$ is a constant function.
\end{solution}

\section{Exercise 4.58}

For biholomorphic functions $f, g : D \to D$ with $f(0) = g(0)$ and
$f'(0) = g'(0)$, show that $f = g$.

\begin{solution}
  Consider $h = g^{-1} \circ f : D \to D$. Then $h$ is biholomorphic
  and $h(0) = g^{-1}(f(0)) = g^{-1}(g(0)) = 0$.

  By the inverse function theorem, $(g^{-1})'(g(0)) = 1/g'(0)$, so
  \[
    h'(0) = (g^{-1})'(f(0)) \cdot f'(0)
    = \frac{f'(0)}{g'(0)} = 1.
  \]

  By Lemma~4.56 (Schwarz lemma), $|h'(0)| \le 1$. Since
  $|h'(0)| = 1$, the equality case yields $h(z) = \zeta z$ for some
  $\zeta$ with $|\zeta| = 1$. As $h'(0) = \zeta = 1$, we conclude
  $h(z) = z$, i.e., $g^{-1} \circ f = \mathrm{id}$, hence $f = g$.
\end{solution}

\section{Exercise 4.59 (Study 1913)}

For a biholomorphic function $f : D \to \Omega$, show that if $\Omega$
is convex, then for any $r \in (0, 1)$, the image
$\Omega_r = f(U_r(0))$ is also convex.

\textit{Hint:} For two points $|a| \le |b| < r$ and $t \in [0, 1]$,
consider the map
$g_t(z) = (1-t)\,f\!\left(\frac{a}{b}z\right) + t\,f(z)$.

\begin{solution}
  Take $w_1 = f(a)$ and $w_2 = f(b)$ in $\Omega_r$ with
  $|a| \le |b| < r$. We show $(1-t)w_1 + tw_2 \in \Omega_r$ for all
  $t \in [0,1]$.

  Define
  $g_t(z) = (1-t)\,f\!\left(\tfrac{a}{b}z\right) + t\,f(z)$
  for $z \in D$. Since $|a/b| \le 1$, for any $z \in D$ we have
  $|\tfrac{a}{b}z| < 1$, so both $f(\tfrac{a}{b}z)$ and $f(z)$ lie
  in $\Omega$. By convexity of $\Omega$, $g_t(z) \in \Omega$ for all
  $z \in D$.

  Now define $h_t = f^{-1} \circ g_t : D \to D$. Then $h_t$ is
  analytic and
  $h_t(0) = f^{-1}(g_t(0)) = f^{-1}(f(0)) = 0$.
  By Lemma~4.56 (Schwarz lemma),
  $|h_t(z)| \le |z|$ for all $z \in D$.
  In particular,
  \[
    |h_t(b)| \le |b| < r.
  \]
  Since
  $h_t(b) = f^{-1}\!\bigl((1-t)f(a) + tf(b)\bigr)$
  and $|h_t(b)| < r$, the point $h_t(b)$ lies in $U_r(0)$, so
  $(1-t)f(a) + tf(b) \in f(U_r(0)) = \Omega_r$.
\end{solution}

\section{Exercise 4.63}

Consider an analytic function $f : D \to D$ that is not the identity
function.
\begin{enumerate}[(a)]
  \item Show that $f$ has at most one fixed point.
  \item Give an example of $f$ having no fixed points.
\end{enumerate}

\begin{solution}
  (a) Suppose $f$ has two distinct fixed points $a, b \in D$.
  Consider the map $g = \varphi_a \circ f \circ \varphi_a^{-1} : D \to
  D$, where $\varphi_a$ is as defined in~(4.18). Then
  \[
    g(0) = \varphi_a(f(\varphi_a^{-1}(0))) = \varphi_a(f(a))
    = \varphi_a(a) = 0.
  \]
  By Lemma~4.56, $|g(z)| \le |z|$ for all $z \in D$. Since $b \ne a$,
  the point $\varphi_a(b) \ne 0$ is also a fixed point of $g$:
  \[
    g(\varphi_a(b)) = \varphi_a(f(b)) = \varphi_a(b).
  \]
  Hence $|g(\varphi_a(b))| = |\varphi_a(b)|$, which is equality in the
  Schwarz lemma. Therefore $g(z) = \zeta z$ for some $|\zeta| = 1$, and
  $\zeta \varphi_a(b) = \varphi_a(b)$ gives $\zeta = 1$, so
  $g = \mathrm{id}$. But then $f = \varphi_a^{-1} \circ \varphi_a =
  \mathrm{id}$, a contradiction.

  \medskip
  (b) Consider $f(z) = -\varphi_a(z) = \frac{a - z}{\bar{a}z - 1}$
  for any $a \in D \setminus \{0\}$.
  This is an automorphism of $D$ by Theorem~4.65 (with $\theta = \pi$),
  so $f : D \to D$ is analytic and not the identity.
  A fixed point satisfies $f(z) = z$, i.e.,
  $\frac{a - z}{\bar{a}z - 1} = z$, which gives
  $a - z = \bar{a}z^2 - z$, hence $\bar{a}z^2 = a$,
  i.e., $z^2 = a/\bar{a}$.
  Therefore $|z|^2 = |a|/|\bar{a}| = 1$, so $|z| = 1$ for both roots,
  and neither lies in $D$.
  Hence $f$ has no fixed points in $D$.
\end{solution}

\section{Exercise 4.64}

For an analytic function $f : D \to D$, show that if there exists a
point $a \in D \setminus \{0\}$ such that $f(0) = a$ and $f(a) = 0$,
then $f = \varphi_a$, where $\varphi_a$ is as defined in~(4.18).

\begin{solution}
  Consider $h = \varphi_a \circ f : D \to D$, which is analytic.
  We compute
  \[
    h(0) = \varphi_a(f(0)) = \varphi_a(a) = 0.
  \]
  By Lemma~4.56, $|h(z)| \le |z|$ for all $z \in D$. On the other
  hand, using Lemma~4.62(b):
  \[
    h(a) = \varphi_a(f(a)) = \varphi_a(0) = a.
  \]
  So $|h(a)| = |a|$, which is equality in the Schwarz lemma. Therefore
  $h(z) = \zeta z$ for some $|\zeta| = 1$, and $\zeta = h'(0)$.
  Since $\zeta a = h(a) = a$ and $a \ne 0$, we get $\zeta = 1$.
  Hence $h(z) = z$, i.e., $\varphi_a \circ f = \mathrm{id}$.
  By Lemma~4.62(c), $\varphi_a^{-1} = \varphi_a$, so
  $f = \varphi_a^{-1} = \varphi_a$.
\end{solution}

\section{Exercise 4.67}

Prove Theorem~4.66 (Schwarz--Pick).

\begin{solution}
  Fix $z_0 \in D$ and set $w_0 = f(z_0)$. Define the composition
  \[
    g = \varphi_{w_0} \circ f \circ \varphi_{z_0}^{-1} : D \to D,
  \]
  where $\varphi_a$ is as in~(4.18). Since $\varphi_{z_0}^{-1}(0) =
  \varphi_{z_0}(0) = z_0$ by Lemma~4.62(b,c), we have
  \[
    g(0) = \varphi_{w_0}(f(z_0)) = \varphi_{w_0}(w_0) = 0.
  \]

  By Lemma~4.56, $|g(z)| \le |z|$ and $|g'(0)| \le 1$. We compute
  $g'(0)$ by the chain rule. Since
  $\varphi_a(z) = \frac{z - a}{\bar{a}z - 1}$, we have
  $\varphi_a'(z) = \frac{|a|^2 - 1}{(\bar{a}z - 1)^2}$, so
  \[
    \varphi_{w_0}'(w_0) = \frac{|w_0|^2 - 1}{(|w_0|^2 - 1)^2}
    = \frac{1}{|w_0|^2 - 1} = \frac{-1}{1 - |w_0|^2},
    \qquad
    (\varphi_{z_0}^{-1})'(0) = \varphi_{z_0}'(0) = |z_0|^2 - 1.
  \]
  Therefore
  \[
    g'(0) = \varphi_{w_0}'(w_0) \cdot f'(z_0) \cdot
    (\varphi_{z_0}^{-1})'(0)
    = \frac{f'(z_0)\,(1 - |z_0|^2)}{1 - |w_0|^2}.
  \]
  The condition $|g'(0)| \le 1$ gives
  \[
    |f'(z_0)| \le \frac{1 - |f(z_0)|^2}{1 - |z_0|^2}.
  \]

  Equality holds iff $g(z) = \zeta z$ for some $|\zeta| = 1$, i.e.,
  $\varphi_{w_0} \circ f \circ \varphi_{z_0}^{-1} = \zeta \,\mathrm{id}$,
  which rearranges to
  $f = \varphi_{w_0}^{-1} \circ (\zeta\,\mathrm{id}) \circ \varphi_{z_0}$.
  This is a composition of automorphisms of $D$ (a rotation and two
  M\"obius maps), hence $f \in \mathrm{Aut}(D)$.

  Conversely, if $f \in \mathrm{Aut}(D)$, then
  $f(z) = e^{i\theta}\varphi_a(z)$ by Theorem~4.65, and one verifies
  directly that equality holds everywhere.
\end{solution}

\section{Exercise 4.74}

The ray in the proof of Theorem~4.73 starts from $f(x)$. Can we start
the ray from $x$ and still get $g$ as a retraction? Why?

\begin{solution}
  No. If we define $g(x)$ as the intersection of the ray starting from
  $x$ in the direction of $f(x)$ with $S^{n-1}$, then for a boundary
  point $p \in S^{n-1}$, the ray from $p$ towards $f(p) \in D^n$
  points inward. This ray intersects $S^{n-1}$ at a point different
  from $p$ (on the opposite side of $D^n$), so $g(p) \ne p$.

  For a concrete example, take $n = 1$, so $D^1 = [-1,1]$ and
  $S^0 = \{-1, 1\}$. Let $f(x) = -x$ (which has no fixed point in
  $D^1$). For $x = 1$, the ray from $1$ towards $f(1) = -1$ gives
  $g(1) = -1 \ne 1$, violating the retraction condition
  $g|_{S^0} = \mathrm{id}_{S^0}$.
\end{solution}

\section{Exercise 4.82}

Prove Corollary~4.81 by similar arguments as in the proof of
Theorem~4.79.

\begin{solution}
  Define the inversion $\Phi(z) = r^2/\bar{z}$, which maps $D_+$ to
  $D_-$ and vice versa, and fixes each point of $D_0$ (since
  $|z| = r$ implies $\bar{z} = r^2/z$, so $\Phi(z) = z$).  Define
  \[
    F(z) =
    \begin{cases}
      f(z) & \text{if } z \in D_+ \cup D_0, \\[4pt]
      \dfrac{\rho^2}{\overline{f(\Phi(z))}} & \text{if } z \in D_-.
    \end{cases}
  \]

  \textbf{Continuity on $D_0$.}
  For $z_0 \in D_0$, $\Phi(z_0) = z_0$ and $|f(z_0)| = \rho$, so
  $\overline{f(z_0)} = \rho^2/f(z_0)$, giving
  $F(z_0) = \rho^2/\overline{f(z_0)} = f(z_0)$.

  \textbf{Analyticity on $D_-$.}
  The map $z \mapsto f(\Phi(z)) = f(r^2/\bar{z})$ is anti-holomorphic
  (composition of holomorphic $f$ with anti-holomorphic $\Phi$).
  Conjugation turns anti-holomorphic into holomorphic, so
  $\overline{f(\Phi(z))}$ is holomorphic on $D_-$.
  Hence $F(z) = \rho^2/\overline{f(\Phi(z))}$ is analytic on $D_-$
  wherever the denominator is nonzero; near $D_0$ this is guaranteed
  since $|f| = \rho > 0$ on $D_0$.

  \textbf{Analyticity on $D$.}
  By the same Morera-type argument as in the proof of
  Theorem~4.79---decomposing arbitrary triangles by their intersection
  with $\partial U_r(0)$ and applying the Cauchy integral
  theorem---$F$ is analytic on all of $D$.
\end{solution}

\section{Exercise 4.83}

Show that an analytic function $f : D \to \C$ that is continuous on
$\overline{D}$ must be a constant if
\[
  f(z) \ne 0 \text{ if } z \in D, \quad
  |f(z)| = 1 \text{ if } z \in \partial D.
\]

\begin{solution}
  Since $f \ne 0$ on $D$, the function $1/f$ is analytic on $D$.
  On $\partial D$, $|f(z)| = 1$.

  By the maximum modulus principle (Theorem~4.48),
  $|f(z)| \le \max_{\partial D}|f| = 1$ for all $z \in D$.
  Applying the same principle to $1/f$:
  $|1/f(z)| \le \max_{\partial D}|1/f| = 1$, i.e.,
  $|f(z)| \ge 1$ for all $z \in D$.

  Combining these, $|f(z)| = 1$ for all $z \in D$.
  Since $|f|$ is constant, $f(D) \subseteq \partial D$, which is not
  open in $\C$. By the open mapping theorem (Theorem~4.45), $f$ must
  be constant.
\end{solution}

\section{Exercise 4.86}

Show that the real part and the imaginary part of an analytic function
is real analytic in the sense of Definition~4.85.

\begin{solution}
  Let $f = u + iv$ be analytic on a domain $\Omega \subseteq \C$,
  identified with an open subset of $\R^2$. Fix $(x_0, y_0) \in \Omega$
  and set $z_0 = x_0 + iy_0$.

  By Theorem~4.1, $f$ has a convergent power series expansion
  \[
    f(z) = \sum_{n=0}^{\infty} a_n (z - z_0)^n,
    \qquad a_n \in \C,
  \]
  valid in some disk $U_r(z_0)$. Expanding $(z - z_0)^n$ by the
  binomial theorem with $z - z_0 = (x - x_0) + i(y - y_0)$:
  \[
    f(z) = \sum_{n=0}^{\infty} a_n
    \sum_{k=0}^{n} \binom{n}{k}
    (x - x_0)^k \cdot i^{n-k}(y - y_0)^{n-k}.
  \]
  Since the series converges absolutely, we may rearrange. Setting
  $m = k$ and $j = n - k$:
  \[
    u(x,y) = \sum_{m=0}^{\infty} \sum_{j=0}^{\infty}
    \Re\!\left(a_{m+j}\,\binom{m+j}{m}\,i^{\,j}\right)
    (x - x_0)^m (y - y_0)^j.
  \]
  This is a convergent power series in two variables centered at
  $(x_0, y_0)$, so $u$ is real analytic by Definition~4.85.
  The same argument with $\Im$ in place of $\Re$ shows that $v$ is
  real analytic.
\end{solution}

\end{document}
